\documentclass[10pt,a4paper]{amsart}
\usepackage[margin=19mm]{geometry}
\usepackage{amsmath,amssymb,mathtools}
\usepackage[scr=rsfs,cal=euler]{mathalfa}
\usepackage{xfrac}
\usepackage[unicode,hidelinks,bookmarks=false]{hyperref}
\newcommand{\ie}{i.e.\ }
\newcommand{\cf}{cf.\ }
\newcommand{\resp}{respectively\ }
\newcommand{\Subsection}{\subsection}
\providecommand{\nobreakdash}{\nobreak\hbox{-}}
\setlength{\emergencystretch}{2em}
\multlinegap0pt
\usepackage{bm}
\usepackage{bbm}
\usepackage{dsfont}
\usepackage{xspace}
\usepackage{cancel}
\usepackage{mathtools}
\usepackage{amsthm}
\makeatletter
\@ifundefined{theorem}{\newtheorem{theorem}{Theorem}}{}
\makeatother
\title[Relative Invariants from Moving Frames on an Extended Manifo]{Relative Invariants from Moving Frames on an Extended Manifold}
\author{Leonid Bedratyuk}
\date{Source: arXiv:2601.06660v1 (CC BY 4.0)}
\begin{document}
\maketitle

\noindent\emph{Source and licence.} Relative Invariants from Moving Frames on an Extended Manifold. Version arXiv:2601.06660v1 (CC BY 4.0), distributed under the Creative Commons Attribution 4.0 International licence, as recorded in the arXiv licence metadata for this exact version and shown on the abstract page https://arxiv.org/abs/2601.06660v1. Full rights notice: licenses/arxiv-2601.06660-CC-BY-4.0.txt.\par\smallskip

\noindent\textit{Editorial scope.} Counted unit: the Theorem whose source label is \texttt{twisted-action} in the article (source0.tex lines 548--555, source label \texttt{twisted-action}), delivered together with the article's own complete original proof of it (source0.tex lines 557--633, one \texttt{proof} environment of 77 lines). No mathematics is written, repaired or re-derived: statement and proof are the article's own text line for line. Required context retained uncounted: twisted-xn1 (lines 538--544). Every definition, display and label of those lines is retained unchanged. Omitted from this package and not redistributed: 1--537, 545--547, 556--556, 634--1311. Nothing inside a retained excerpt is omitted or altered: every retained body is delivered exactly as the article's lines are written, so no omission receipt of any kind accompanies this package. Citations: the retained excerpts carry no citation command at all, so no bibliography entry of the article is transcribed and no reference is redistributed. Reference adaptation: none was needed and none was made; every reference inside the delivered unit resolves inside it, so no unresolved reference or citation is produced. Printed slips kept literal: no printed word, symbol, hypothesis or number of the counted statement or of its proof was altered. Layout and class: the article's own class is \texttt{amsart} with packages including latexsym, amsmath, amscd, amssymb, amsthm, longtable, fontenc, inputenc, babel, mathtools, cite; the wrapper is the lane's pinned \texttt{amsart} editorial wrapper at 10pt on A4, which restates the theorem-like environments the retained text needs and adds the article's own macro definitions that the retained text needs (none), with extra wrapper packages (bm, bbm, dsfont, xspace, cancel, mathtools). The delivered numbering is the unit's own. Assets: no retained span contains a figure environment, a graphics-inclusion command, a PDF-page inclusion, a table or a TikZ picture; no image file is copied into this package, the package declares no asset dependency and no image file is redistributed with it. Licence: version arXiv:2601.06660v1 of this article is distributed under the Creative Commons Attribution 4.0 International licence, as recorded in the arXiv licence metadata for this exact version and shown on the abstract page https://arxiv.org/abs/2601.06660v1. A full rights notice is delivered at licenses/arxiv-2601.06660-CC-BY-4.0.txt. Characters: every retained excerpt is pure ASCII, so the delivered engine, XeLaTeX with its default Latin Modern text font, reports no missing character.
\begin{equation}\label{twisted-xn1}
g \cdot \boldsymbol{\hat x}
\;=\;
\bigl(g \cdot \boldsymbol x,\ x_{m+1} \cdot \mu(g,\boldsymbol x)\bigr),
\quad
\boldsymbol{\hat x} = (x_1, \dots, x_m, x_{m+1}) \in \widehat{\mathcal M}.
\end{equation}

\begin{theorem}\label{twisted-action}
The following statements hold:
\begin{enumerate}\itemsep=2pt
\item[(\textit{i})] The map \eqref{twisted-xn1} defines a smooth action of $G$ on $\widehat{\mathcal M}$.
\item[(\textit{ii})] The projection $\pi$ is a $G$-equivariant map.
\item[(\textit{iii})] If the action of $G$ on $\mathcal M$ is free and regular near a point $\boldsymbol x$, then the induced action of $G$ on $\widehat{\mathcal M}$ is also free and regular near any point in $\pi^{-1}(\boldsymbol x)$.
\end{enumerate}
\end{theorem}

\begin{proof}
(\textit{i}) For the identity $e \in G$, we compute:
$$
e \cdot \boldsymbol{\hat x}
= (e \cdot \boldsymbol x,\ x_{m+1} \mu(e,\boldsymbol x))
= (\boldsymbol x,\ x_{m+1}) = \boldsymbol{\hat x},
$$
since $\mu(e, \boldsymbol x) = 1$.

For any $g,h \in G$, using the cocycle identity for $\mu$, we have
\begin{align*}
g \cdot (h \cdot \boldsymbol{\hat x})
&= g \cdot \bigl(h \cdot \boldsymbol x,\ x_{m+1} \mu(h,\boldsymbol x)\bigr) \\
&= \bigl(gh \cdot \boldsymbol x,\ x_{m+1} \mu(h,\boldsymbol x)\, \mu(g, h \cdot \boldsymbol x)\bigr) \\
&= \bigl(gh \cdot \boldsymbol x,\ x_{m+1} \mu(gh,\boldsymbol x)\bigr)
= (gh) \cdot \boldsymbol{\hat x}.
\end{align*}
Smoothness follows from the smoothness of the action on $\mathcal M$ and of the function $\mu$.

\medskip

(\textit{ii}) To verify equivariance of $\pi$:
$$
\pi(g \cdot \boldsymbol{\hat x}) = \pi(g \cdot \boldsymbol x,\ x_{m+1} \mu(g,\boldsymbol x)) = g \cdot \boldsymbol x = g \cdot \pi(\boldsymbol{\hat x}).
$$

\medskip

(\textit{iii}) Freeness: if the stabilizer $G_{\boldsymbol x} = \{ g \in G \mid g \cdot \boldsymbol x = \boldsymbol x \}$ is trivial, then so is
\[
G_{\boldsymbol{\hat x}} = \{ g \in G \mid g \cdot \boldsymbol x = \boldsymbol x \ \text{and} \ x_{m+1} \mu(g, \boldsymbol x) = x_{m+1} \}
= \{ g \in G_{\boldsymbol x} \mid \mu(g,\boldsymbol x) = 1 \} = \{e\}.
\]

\medskip

Semi-regularity: The orbit of a point $\boldsymbol{\hat x} \in \widehat{\mathcal M}$ takes the form
\[
\mathcal O_{\widehat{\mathcal M}}(\boldsymbol{\hat x})
= \{ (g \cdot \boldsymbol x,\ x_{m+1} \mu(g,\boldsymbol x)) \mid g \in G \}.
\]
Since $\pi$ is $G$-equivariant, it maps orbits to orbits, and the restriction
\[
\pi : \mathcal O_{\widehat{\mathcal M}}(\boldsymbol{\hat x}) \to \mathcal O_{\mathcal M}(\boldsymbol x)
\]
is bijective. Indeed, for each $\boldsymbol y \in \mathcal O_{\mathcal M}(\boldsymbol x)$, there exists a unique $g_{\boldsymbol y} \in G$ such that $\boldsymbol y = g_{\boldsymbol y} \cdot \boldsymbol x$; then, using the cocycle identity for $\mu$ and the equivariance of $\pi$, the inverse map to $\pi|_{\mathcal O_{\widehat{\mathcal M}}(\boldsymbol{\hat x})}$ is given by
\[
\boldsymbol y \mapsto \bigl(\boldsymbol y,\ x_{m+1} \mu(g_{\boldsymbol y}, \boldsymbol x)\bigr).
\]
Moreover, this inverse is smooth, hence $\pi$ restricts to a diffeomorphism between orbits, and
\[
\dim \mathcal O_{\widehat{\mathcal M}}(\boldsymbol{\hat x}) = \dim \mathcal O_{\mathcal M}(\boldsymbol x),
\]
so semi-regularity is preserved.

\medskip

Regularity: Let the action of $G$ be free and regular in a neighborhood $U \subset \mathcal M$ of $\boldsymbol x$, so that the intersection $U \cap \mathcal O_{\mathcal M}(\boldsymbol x)$ is connected.
Choose $\varepsilon > 0$ such that the interval $(x_{m+1} - \varepsilon,\ x_{m+1} + \varepsilon) \subset \mathbb{R}^\times$ and set
\[
\widehat U = U \times (x_{m+1} - \varepsilon,\ x_{m+1} + \varepsilon) \subset \widehat{\mathcal M}.
\]
Then $\pi$ maps $\widehat U \cap \mathcal O_{\widehat{\mathcal M}}(\boldsymbol{\hat x})$ onto $U \cap \mathcal O_{\mathcal M}(\boldsymbol x)$.
Since the action is free, for each $\boldsymbol y \in U \cap \mathcal O_{\mathcal M}(\boldsymbol x)$ there exists a unique $g_{\boldsymbol y} \in G$ such that $\boldsymbol y = g_{\boldsymbol y} \cdot \boldsymbol x$. Due to the regularity of the orbit map $G \to \mathcal O_{\mathcal M}(\boldsymbol x)$, the map $\boldsymbol y \mapsto g_{\boldsymbol y}$ can be chosen locally smoothly. Define the map
\[
\Psi : U \cap \mathcal O_{\mathcal M}(\boldsymbol x) \longrightarrow \widehat U \cap \mathcal O_{\widehat{\mathcal M}}(\boldsymbol{\hat x}), \qquad
\Psi(\boldsymbol y) = \bigl(\boldsymbol y,\ x_{m+1} \mu(g_{\boldsymbol y}, \boldsymbol x)\bigr).
\]
This map is well-defined and smooth, and its inverse is given by $\pi$, since
\[
\pi(\Psi(\boldsymbol y)) = \boldsymbol y, \qquad \Psi(\pi(\boldsymbol{\hat y})) = \boldsymbol{\hat y}
\]
for all $\boldsymbol{\hat y} \in \widehat U \cap \mathcal O_{\widehat{\mathcal M}}(\boldsymbol{\hat x})$.
Hence $\Psi$ is a diffeomorphism.

Thus, $\widehat U \cap \mathcal O_{\widehat{\mathcal M}}(\boldsymbol{\hat x})$ is diffeomorphic to a connected subset, and the group action on $\widehat{\mathcal M}$ is regular near $\boldsymbol{\hat x}$, as claimed.
\end{proof}

\end{document}
